% fixture: babel/portuguese-basic
% tier: T3
% engine: document
% covers: \usepackage[portuguese]{babel} (portuges.ldf): its captions (Resumo, \'Indice in the contents' heading and running heads, Figura, Tabela, Demonstra\c{c}\~ao, Refer\^encias), \today, the portuguese patterns (hyphenmins 2 3, no \frenchspacing), its " shorthands ("" "- and "< "> as babel's OT1 guillemets), babel's OT1 umlauts
% issue: #63
\documentclass{article}
\usepackage{amsthm}
\usepackage[portuguese]{babel}
\pagestyle{headings}
\title{Título do artigo}
\author{João Silva}
\begin{document}
\maketitle
\begin{abstract}
Um resumo curto. Fim.
\end{abstract}
\tableofcontents
\section{Introdução}
A teoria da informação estuda a quantificação, o armazenamento e a comunicação da informação digital.
Palavras com hífen: guarda"-chuva e couve"-flor, sem""quebra, aspas "<assim"> e na\"{\i}ve, \"a, ü.
\begin{figure}[h]
\caption{Uma figura}
\end{figure}
\begin{table}[h]
\caption{Uma tabela}
\end{table}
\begin{proof}
Uma demonstração.
\end{proof}
\begin{minipage}{60pt}
Transformações estabelecidas e armazenamento de comunicações eletrônicas.
\end{minipage}
\begin{thebibliography}{9}
\bibitem{a} Um item.
\end{thebibliography}
\end{document}
